\documentclass[11pt]{article}
\usepackage[a4paper, top=18mm, bottom=18mm, left=20mm, right=20mm]{geometry}
\usepackage[T2A]{fontenc}
\usepackage[utf8]{inputenc}
\usepackage[ukrainian]{babel}
\usepackage{graphicx}
\usepackage[export]{adjustbox}
\usepackage{amsmath}
\usepackage{amssymb}
\usepackage{microtype}
\graphicspath{{/app/Quiz/Scripts/2023/ProgAll/15BinTree/}}
\setlength{\parindent}{0pt}
\setlength{\parskip}{6pt}
\emergencystretch=3em
\pagestyle{empty}
\begin{document}
%begin
{\centering\Large\bfseries Контрольна робота\par}
\medskip
\noindent\rule{\linewidth}{0.5pt}
\smallskip
\noindent\textbf{Варіант \No\,1}\hfill \textit{ПІБ:}\,\underline{\hspace{60mm}}
\vspace{4mm}

%beginitem
1. %goodpath:Scripts/2018/Mod2018_1/03-good.txt
\textbf{Відмітити все, що є коректним оператором С++:}
( 1 ) cout<<'\n';

( 2 ) int n,n;

( 3 ) if ('x'><'y') a=6;

( 4 ) {a=5;}
\vspace{5mm}
%enditem
%beginitem
2. \textbf{Що буде виведено за виконання фрагменту коду?}

%code
-3**True+-2
%code
\vspace{5mm}
%enditem
%beginitem
3. \textbf{Що буде виведено за виконання фрагменту коду?}

%code
print(2 == True <= 8 != 2.0)
%code
\vspace{5mm}
%enditem
%beginitem
4. %goodpath:Scripts/2019/Mod2019_1/Lex/01-good.txt
\textbf{Відмітити все, що є однією правильною лексемою:}
( 1 ) {

( 2 ) <+

( 3 ) =<

( 4 ) 0X2b
\vspace{5mm}
%enditem
%beginitem
5. \textbf{Що буде виведено за виконання фрагменту коду?}

a,b,c=0,6,9
def h(a):
global c    a=5
    b*=2
    c=1
    return a+b+c

a,b,c=7,8,3
print(h(a),a,b,c)
\vspace{5mm}
%enditem
%beginitem
6. \textbf{Що буде виведено за виконання фрагменту коду?}

%code
cout << (8 % 8 % 8 % 8);
%code
\vspace{5mm}
%enditem
%beginitem
7. 

\begin{center}
	\adjustimage{max size={0.8\linewidth}{0.8\paperheight}}
	{../15BinTree/trees_exp/86.png}
\end{center}
Указати послідовність міток вершин дерева, зображеного на рис., що утвориться в результаті обходу дерева в прямому порядку. Мітки записувати через пробіл.\par Алгоритм починає роботу з кореня (на рис. це верхня вершина).
%answer:86 прямому
\vspace{5mm}
%enditem
%beginitem
8. %goodpath:Scripts/2019/Mod2019_1/Lex/03-good.txt
\textbf{Відмітити все, що є коректним оператором С++:}
( 1 ) F"1={x:5.2f};"

( 2 ) 'a'=3

( 3 ) True=1

( 4 ) x=float(4)
\vspace{5mm}
%enditem
%beginitem
9. \textbf{Що буде виведено за виконання фрагменту коду?}
%code
a, b, c = 6, 9, 8
def h(a, b, c):
    print(a, b, c, end=" ")

h(0, 5)
print(a, b, c)
%code
\vspace{5mm}
%enditem
%beginitem
10. \textbf{Що буде виведено за виконання фрагменту коду?}

print('R', end=' ')
a=[10,11,12,13,14]
a.append('13')
print(a)

\vspace{5mm}
%enditem










%end
\newpage

\newpage

%begin
{\centering\Large\bfseries Контрольна робота\par}
\medskip
\noindent\rule{\linewidth}{0.5pt}
\smallskip
\noindent\textbf{Варіант \No\,2}\hfill \textit{ПІБ:}\,\underline{\hspace{60mm}}
\vspace{4mm}

%beginitem
1. \textbf{Що буде виведено за виконання фрагменту коду?}

a,b,c=0,6,9
def h(a):
global c    a=5
    b*=2
    c=1
    return a+b+c

a,b,c=7,8,3
print(h(a),a,b,c)
\vspace{5mm}
%enditem
%beginitem
2. \textbf{Що буде виведено за виконання фрагменту коду?}

%code
print('R', end=' ')
a = ('a','b','c','d','e',0,1,2,3,4)
print(a[4])
%code

\vspace{5mm}
%enditem
%beginitem
3. %goodpath:Scripts/2018/Mod2018_1/03-good.txt
\textbf{Відмітити все, що є коректним оператором С++:}
( 1 ) cout<<'\n';

( 2 ) int n,n;

( 3 ) if ('x'><'y') a=6;

( 4 ) {a=5;}
\vspace{5mm}
%enditem
%beginitem
4. \textbf{Що буде виведено за виконання фрагменту коду?}

%code
cout << (!2<true and !8!= 2.0 and 6==8.0);
%code
\vspace{5mm}
%enditem
%beginitem
5. \textbf{Що буде виведено за виконання фрагменту коду?}
%code

a,b,c=6,9,8
def h(a,b,c):
    print(a,b,c,end="")

h(0,5)
print(a,b,c)

%code
\vspace{5mm}
%enditem
%beginitem
6. \textbf{Що буде виведено за виконання фрагменту коду?}

x,y,z=0,5,4
x,z,y,y=6,z,y,9
print(x,y,z)
\vspace{5mm}
%enditem
%beginitem
7. %goodpath:Scripts/2019/Mod2019_1/Lex/03-good.txt
\textbf{Відмітити все, що є коректним оператором С++:}
( 1 ) F"1={x:5.2f};"

( 2 ) 'a'=3

( 3 ) True=1

( 4 ) x=float(4)
\vspace{5mm}
%enditem
%beginitem
8. %goodpath:Scripts/2018/Mod2018_1/01-good.txt
\textbf{Відмітити все, що є однією правильною лексемою:}
( 1 ) {

( 2 ) <+

( 3 ) <>

( 4 ) =
\vspace{5mm}
%enditem
%beginitem
9. \textbf{Що буде виведено за виконання фрагменту коду?}

%code
print(-3**True+-2)
%code
\vspace{5mm}
%enditem
%beginitem
10. \textbf{Що буде виведено за виконання фрагменту коду?}

%code
#include <iostream>
using namespace std;

int h(int a){
    int u = 96;
    if (a < -1) 
        return 0;
    if (a != 2)
         u = 6;
    else
         return 9;
    return u;
}

int main(){
    cout << h(2);
    return 0;
}
%code

\vspace{5mm}
%enditem










%end
\newpage

\newpage

%begin
{\centering\Large\bfseries Контрольна робота\par}
\medskip
\noindent\rule{\linewidth}{0.5pt}
\smallskip
\noindent\textbf{Варіант \No\,3}\hfill \textit{ПІБ:}\,\underline{\hspace{60mm}}
\vspace{4mm}

%beginitem
1. \textbf{Що буде виведено за виконання фрагменту коду?}

%code
cout << 8 % 8 % 8 % 8;
%code
\vspace{5mm}
%enditem
%beginitem
2. %goodpath:Scripts/2019/Mod2019_1/Lex/01-good.txt
\textbf{Відмітити все, що є однією правильною лексемою:}
( 1 ) {

( 2 ) <+

( 3 ) =<

( 4 ) 0X2b
\vspace{5mm}
%enditem
%beginitem
3. \textbf{Що буде виведено за виконання фрагменту коду?}
try:
    res = -3**True+-2
    if isinstance(res, bool):
        print(f'bool;{res}')
    elif isinstance(res, int):
        print(f'int;{res}')
    elif isinstance(res, float):
        print(f'float;{res:.2f}')
    else:
        print('???')
except:
    print('Z')
\vspace{5mm}
%enditem
%beginitem
4. \textbf{Що буде виведено за виконання фрагменту коду?}
try:
    res = 7/7*7*7
    if isinstance(res, bool):
        print(f'bool;{res}')
    elif isinstance(res, int):
        print(f'int;{res}')
    elif isinstance(res, float):
        print(f'float;{res:.2f}')
    else:
        print('???')
except:
    print('Z')
\vspace{5mm}
%enditem
%beginitem
5. 
Вкажіть значення та тип виразу:
5<"True"
\vspace{5mm}
%enditem
%beginitem
6. %goodpath:Scripts/2019/Mod2019_1/Lex/03-good.txt
\textbf{Відмітити все, що є коректним оператором С++:}
( 1 ) F"1={x:5.2f};"

( 2 ) 'a'=3

( 3 ) True=1

( 4 ) x=float(4)
\vspace{5mm}
%enditem
%beginitem
7. \textbf{Що буде виведено за виконання фрагменту коду?}
try:
    for a in range(-4, -4, 3):
        if a<3:
            break
        print(a, end=';')
        if a<3:
            break
    else:
        print('end',end=';')
    print(a)
except:
    print('Z')
\vspace{5mm}
%enditem
%beginitem
8. \textbf{Що буде виведено за виконання фрагменту коду?}
%code
#include <iostream>
using namespace std;

int f(int &x, int &y){
    x = 5;
    y+= 8;
    return x;
}

int main(){
    int a = 6, b = 4;
    a = f(a, a);
    cout << a << ":" << b <<':';
    {
        int a = 7, b = 1;
        cout << ((a<7) && ((b+=1) < 7))
             << a << ":" << b << ":";
    }
    {
        inta = 9;
        cout << a << ':';
    }
    cout << a << ":" << b << endl;
    return 0;
}
%code

\vspace{5mm}
%enditem
%beginitem
9. %goodpath:Scripts/2018/Mod2018_1/01-good.txt
\textbf{Відмітити все, що є однією правильною лексемою:}
( 1 ) {

( 2 ) <+

( 3 ) <>

( 4 ) =
\vspace{5mm}
%enditem
%beginitem
10. \textbf{Що буде виведено за виконання фрагменту коду?}

%code
print('R', end=' ')
i=0
while i<9:
    i+=1
print(i)
%code

\vspace{5mm}
%enditem










%end
\newpage

\newpage

%begin
{\centering\Large\bfseries Контрольна робота\par}
\medskip
\noindent\rule{\linewidth}{0.5pt}
\smallskip
\noindent\textbf{Варіант \No\,4}\hfill \textit{ПІБ:}\,\underline{\hspace{60mm}}
\vspace{4mm}

%beginitem
1. \textbf{Що буде виведено за виконання фрагменту коду?}
a = 0
b = 6
c = 9

def h(a):
global c    a = 5
    b *= 2
    c = 1
    return a + b + c

a = 7
b = 8
c = 3

try:
    print(h(a), a, b, c, sep=';')
except:
    print('Z', a, b, c, sep=';')
\vspace{5mm}
%enditem
%beginitem
2. 
Записати ВИРАЗ, значенням якого є множина, що складається з цілих чисел від 14 до 2011 (включно), що не діляться на 2.\vspace{5mm}
%enditem
%beginitem
3. %goodpath:Scripts/2019/Mod2019_1/Lex/03-good.txt
\textbf{Відмітити все, що є коректним оператором С++:}
( 1 ) F"1={x:5.2f};"

( 2 ) 'a'=3

( 3 ) True=1

( 4 ) x=float(4)
\vspace{5mm}
%enditem
%beginitem
4. 
Записати ВИРАЗ, значенням якого є список, що складається з кубів цілих чисел від 14 до 2011 (включно), що не діляться на жодне з чисел 2, 3, записаних в оберненому порядку.\vspace{5mm}
%enditem
%beginitem
5. %goodpath:Scripts/2018/Mod2018_1/01-good.txt
\textbf{Відмітити все, що є однією правильною лексемою:}
( 1 ) {

( 2 ) <+

( 3 ) <>

( 4 ) =
\vspace{5mm}
%enditem
%beginitem
6. 
Написати програму, що запитує у користувача значення дійсних параметрів $x$ та $y$, обчислює для них значення виразу 
$$\frac{\pi x + 74}{(\arcsin x + 0.5) (y + 98)}$$ 
та повiдомляє введенi данi та результат обчислення.
 Оцінюється правильність та якість коду.

\vspace{5mm}
%enditem
%beginitem
7. %goodpath:Scripts/2019/Mod2019_1/Lex/01-good.txt
\textbf{Відмітити все, що є однією правильною лексемою:}
( 1 ) {

( 2 ) <+

( 3 ) =<

( 4 ) 0X2b
\vspace{5mm}
%enditem
%beginitem
8. \textbf{Що буде виведено за виконання фрагменту коду?}

print('R', end=' ')
for a in range(-4,-4,3):
    if a<3:
        break
        print(a, end=' ')
    if a<3:
        break
    else:
        print('end', end=' ')
    print(a)
\vspace{5mm}
%enditem
%beginitem
9. \textbf{Що буде виведено за виконання фрагменту коду?}
a=[10,11,12,13,14]; a.append('13'); print(a)\vspace{5mm}
%enditem
%beginitem
10. 

\begin{center}
	\adjustimage{max size={0.8\linewidth}{0.8\paperheight}}
	{../15BinTree/trees_exp/86.png}
\end{center}
Указати послідовність міток вершин дерева, зображеного на рис., що утвориться в результаті обходу дерева в прямому порядку. Мітки записувати через пробіл.\par Алгоритм починає роботу з кореня (на рис. це верхня вершина).
%answer:86 прямому
\vspace{5mm}
%enditem










%end
\newpage
\end{document}


